\begin{proposition}[Conductor-cleared jets on a complete-intersection critical curve]\hypertarget{tag-01BF}{}\label{cdp:st:A.5.5}
\StatementRecord{cdp:st:A.5.5}{proved}{01BF}{\SeriesRef{cdp:eq:A.J.4}, \SeriesRef{cdp:eq:A.J.5}, \SeriesRef{cdp:eq:A.J.2}, \SeriesRef{cdp:eq:A.J.3}, \SeriesRef{cdp:eq:A.J.7}, \SeriesRef{cdp:eq:A.J.9}, \SeriesRef{cdp:eq:A.J.8}, \SeriesRef{cdp:eq:A.J.10}}

Let $Y$ be a smooth compact complex threefold with
$H^2(Y,\mathbb Q)=0$, let $S\subset Y$ be a reduced divisor,
and let $D\subset\operatorname{Sing}(S)$ be a reduced
irreducible compact curve which is a local complete intersection
in $Y$. Assume that the generic plane-curve transversal to $D$
in $S$ is an ordinary cusp, an ordinary tacnode, or an ordinary
triple point.
If the normalization $C\to D$ has positive genus, then $D$ is
a smooth elliptic curve. The transversal type is the indicated
cusp, tacnode, or triple point at every point of $D$.

In the cusp case, the intrinsic cubic on the radical of the Hessian
extends to a meromorphic section on $C$ whose pole divisor is bounded
by twice the conductor. The tacnode fourth-order invariant extends
with poles bounded by four times the conductor.
\end{proposition}


\begin{proof}

\begin{enumerate}[label=\arabic*.,leftmargin=*]
\item Let $\nu:C\to D$ be the normalization and $j:C\to Y$
its composite with the inclusion,
$I=\mathcal I_{D/Y}$, and
\[
 N=\bigl(\nu^*(I/I^2)\bigr)^*,\qquad
 M=j^*\mathcal O_Y(S).
\]
These are a rank-two vector bundle and a line bundle on $C$.
Let $\Xi$ be the conductor divisor on $C$, so that
$\nu_*\mathcal O_C(-\Xi)$ is the conductor ideal in
$\mathcal O_D$, and put $c=\deg\Xi\geq0$.
Adjunction for the complete-intersection curve and finite
duality give exact line-bundle comparisons
\[
 \det N=\nu^*\omega_D\otimes j^*K_Y^{-1},\qquad
 \omega_C=\nu^*\omega_D\otimes\mathcal O_C(-\Xi).
                                                        \tag{A.J.1}\hypertarget{tag-005C}{}\label{cdp:eq:A.J.1}
\]
For the second comparison one can work directly in the total
quotient ring: $\operatorname{Hom}_{\mathcal O_D}
(\nu_*\mathcal O_C,\mathcal O_D)$ consists precisely of
elements multiplying the normalization into $\mathcal O_D$,
which is its conductor. Since $D$ is Gorenstein, tensoring
this identity by its invertible dualizing sheaf is finite
duality. Both $M$ and $j^*K_Y$ have degree zero, by the
stated rational cohomology hypothesis. Consequently
\[
 \deg\det N=2g(C)-2+c.                              \tag{A.J.2}\hypertarget{tag-005N}{}\label{cdp:eq:A.J.2}
\]
This is the normal bundle of the complete-intersection ideal,
with its specified normalization pullback.

\item For an ordinary cusp, choose a local defining section
$f$ of $S$, using a local frame of $\mathcal O_Y(S)$.
Since $D\subset\operatorname{Sing}(S)$, its conormal class
maps to zero under
$I/I^2\to\Omega_Y^1|_D$. This homomorphism is injective:
its source is locally free on the reduced curve and the map
is injective at every generic point. Therefore $f\in I^2$.
The associated-graded identity for a regular sequence gives
$I^2/I^3=\operatorname{Sym}^2(I/I^2)$. The class of $f$
is thus a quadratic form
\[
 q\in H^0\bigl(C,\operatorname{Sym}^2N^*\otimes M\bigr).
\]
It has rank one generically. Let $R\subset N$ be its
saturated radical and $Q=N/R$. Both are line bundles.
The quadratic form annihilates $R$ everywhere by holomorphic
continuation, and gives a nonzero section of $M\otimes Q^{-2}$.
In particular
\[
 \deg Q\leq0.                                      \tag{A.J.3}\hypertarget{tag-005Y}{}\label{cdp:eq:A.J.3}
\]

\item On $D_{\rm reg}$, the third-order term on the radical is
intrinsic. Indeed, after writing the quadratic form as
$a\ell^2$, choose holomorphic lifts $\widehat a$ and $\widehat\ell$
and subtract $\widehat a\,\widehat\ell^2$ from $f$.
Restrict the resulting cubic class to $R^3$. Changing these two lifts
adds cubic terms divisible by $\ell$, which vanish on $R^3$.
Changing normal generators has the same effect. Its nonzero
value characterizes the ordinary-cusp transversal among the
rank-one quadratic cases. Denote this cubic by $k$.

\item At a point of $D$, choose two generators $X,Y$ of $I$. After a constant invertible linear
change of these generators, the normalized linear factor of
$q$ has the form $X+hY$ on every normalization branch over
that point. Such a choice exists: the saturated quotient
$N\to Q$ is surjective at each of the finitely many points,
so a constant vector can be chosen outside their finitely
many kernels. Thus $h$ is holomorphic on the normalization,
and
\[
 q=a(X+hY)^2,
\]
where $a=q_{XX}$ is the pullback of a coefficient
in $\mathcal O_D$. Choose a local conductor generator $d$
on the normalization. Both $d$ and $dh$ belong to
$\mathcal O_D$, since the conductor is an ideal in the
normalization contained in that ring. Choose lifts
$\widehat d,\widehat{dh},\widehat a$ to $\mathcal O_Y$,
and set
\[
 G=\widehat d^{\,2} f-
    \widehat a\bigl(\widehat d X+\widehat{dh}Y\bigr)^2.
                                                        \tag{A.J.4}\hypertarget{tag-006C}{}\label{cdp:eq:A.J.4}
\]
The quadratic class of $G$ is zero after normalization.
It is already zero on $D$: the map from the locally free
sheaf $\operatorname{Sym}^2(I/I^2)$ to its normalization
pullback is injective because $D$ is reduced. Hence
$G\in I^3$.

\item Its class in $I^3/I^4=\operatorname{Sym}^3(I/I^2)$,
pulled back to $C$ and evaluated on $R^3$, is holomorphic.
Dividing that class by $d^2$ gives $k$ on the dense regular
locus. To check this equality, there $d$ is a unit and
\eqref{cdp:eq:A.J.4} is exactly the subtraction of the square used in the
definition of the radical cubic. Different choices of lifts
add cubic terms containing its linear factor, and therefore
have zero restriction to $R^3$. Thus these local expressions
are the same meromorphic section wherever they overlap.
Since $d$ has divisor $\Xi$, the resulting bound is
\[
 0\ne k\in H^0\bigl(C,M\otimes R^{-3}
                               \otimes\mathcal O_C(2\Xi)\bigr).
                                                        \tag{A.J.5}\hypertarget{tag-006E}{}\label{cdp:eq:A.J.5}
\]

\item Taking degrees in \eqref{cdp:eq:A.J.5} gives $\deg R\leq2c/3$.
Together with \eqref{cdp:eq:A.J.2}--\eqref{cdp:eq:A.J.3}, this yields
\[
 2g(C)-2+c=\deg R+\deg Q\leq\frac{2c}{3},
 \qquad 2g(C)-2+\frac c3\leq0.                     \tag{A.J.6}\hypertarget{tag-006F}{}\label{cdp:eq:A.J.6}
\]
For $g(C)\geq1$, this forces $g(C)=1$ and $c=0$.
Conductor zero means $\mathcal O_D=\nu_*\mathcal O_C$;
a normal complex curve is smooth. Hence $D=C$ is elliptic.
Now \eqref{cdp:eq:A.J.2}--\eqref{cdp:eq:A.J.5} give $\deg R=\deg Q=0$.
The Hessian coefficient and radical cubic are nonzero
sections of degree-zero line bundles on this elliptic curve,
so neither has a zero. Locally the holomorphic Morse lemma
in the nondegenerate normal direction, with the curve
coordinate as parameter, writes $f=y^2+g(x,z)$.
The conditions just proved give $g=x^3u(x,z)$ with $u$ a
holomorphic unit. Replacing $x$ by a local cube root of $u$
times $x$ gives the local equation $y^2+x^3=0$.

\item In the tacnode case, put $P=\widehat d X+\widehat{dh}Y$ and
$G=\widehat d^{\,2}f-\widehat aP^2\in I^3$.
The generic radical cubic is now zero, so its continuation
$[G]_3|_{R^3}$ is zero everywhere on the normalization.
Since $\ell=X+hY$ is monic in $X$, polynomial division gives
a holomorphic quadratic $C_2$ on the normalization
such that
\[
 [G]_3=\ell C_2.
\]
Each coefficient of $dC_2$ is in the conductor and hence in
$\mathcal O_D$. Lift those coefficients and let $H$ be the
resulting homogeneous quadratic in $X,Y$. Set
\[
 K=\widehat d^{\,2}G-PH.
\]
Its cubic class pulls back to
$d^2\ell C_2-d\ell\cdot dC_2=0$.
The same locally free graded-piece injection used above gives
$K\in I^4$. There is an exact identity in the local ring of
the ambient smooth threefold:
\[
 4\widehat a\widehat d^{\,6}f
 -(2\widehat a\widehat d^{\,2}P+H)^2
 =4\widehat a\widehat d^{\,2}K-H^2.               \tag{A.J.7}\hypertarget{tag-006G}{}\label{cdp:eq:A.J.7}
\]
On the dense smooth locus, complete the square in the
nondegenerate normal direction. If the relevant terms there
are $a\ell^2+b\ell r^2+er^4$, the fourth-order invariant is
$\Delta=ae-b^2/4$. A substitution
$\ell\mapsto\ell+ur^2$ leaves this expression unchanged;
terms divisible by $\ell^2$ and higher changes in the radical
coordinate do not alter it. Thus it is the intrinsic section
of $M^2\otimes Q^{-2}\otimes R^{-4}$ on that dense locus.
It is nonzero precisely for the ordinary-tacnode transversal
among the rank-one Hessian cases with zero radical cubic.

\item Divide \eqref{cdp:eq:A.J.7} by $4\widehat d^{\,6}$ and take the
fourth-order class on $R^4$. Since
$[H]_2=dC_2$, this expresses exactly that same invariant as
\[
 \Delta=
 \frac{a[K]_4|_{R^4}-(C_2|_{R^2})^2/4}{d^4}.
                                                        \tag{A.J.8}\hypertarget{tag-006H}{}\label{cdp:eq:A.J.8}
\]
Its numerator is holomorphic on the normalization.
Different choices agree on the dense regular locus and hence
give one meromorphic section. Therefore
\[
 0\ne\Delta\in H^0\bigl(C,M^2\otimes Q^{-2}
                    \otimes R^{-4}\otimes\mathcal O_C(4\Xi)\bigr).
                                                        \tag{A.J.9}\hypertarget{tag-006I}{}\label{cdp:eq:A.J.9}
\]
Write $q_0=\deg Q\leq0$. Using
$\deg R=2g(C)-2+c-q_0$, the degree of the line bundle
in \eqref{cdp:eq:A.J.9} is
\[
 -2q_0-4\bigl(2g(C)-2+c-q_0\bigr)+4c
       =-8g(C)+8+2q_0.
\]
Its nonzero section and $g(C)\geq1$ force
$g(C)=1$ and $q_0=0$.

\item The section $q$ of $M\otimes Q^{-2}$ is now nowhere zero.
In the local presentation above, $a=q_{XX}$ is therefore a
unit on every normalization branch. Since it belongs to
$\mathcal O_D$, it is a unit in that local ring as well.
The coefficient $h=q_{XY}/a$ consequently belongs to
$\mathcal O_D$; here $q_{XY}$ denotes the coefficient in
the symmetric matrix, so the cross term is $2q_{XY}XY$.
We can thus repeat the entire construction with $d=1$.
Polynomial division by the monic $X+hY$ then takes place
over $\mathcal O_D$ itself; the zero radical cubic descends
by the same injectivity argument. Formula \eqref{cdp:eq:A.J.8} now has
no poles. The resulting nonzero section of
$M^2\otimes Q^{-2}\otimes R^{-4}$ has degree $-4c$,
which forces $c=0$.

\item Thus again $D=C$ is smooth elliptic. Both the nonzero
quadratic coefficient and $\Delta$ are sections of
degree-zero line bundles and are nowhere zero. The parametric
Morse lemma first gives $f=y^2+g(x,z)$, and the vanished
radical cubic and nonzero fourth invariant give
$g=x^4u(x,z)$ with $u$ a unit. Taking a local fourth root
produces the equation $y^2+x^4=0$ at every point of $D$.

\item For an ordinary triple transversal, the quadratic class
vanishes generically, hence vanishes in the locally free
sheaf $I^2/I^3$. Thus $f\in I^3$, and its initial cubic
is a global section
$b\in H^0(C,\operatorname{Sym}^3N^*\otimes M)$.
The discriminant of a binary cubic is homogeneous of degree
four in its coefficients and transforms by the sixth power
of the determinant of a change of variables. Its value is
therefore a section
\[
 0\ne\operatorname{disc}(b)\in
 H^0\bigl(C,M^4\otimes(\det N)^{-6}\bigr).
                                                        \tag{A.J.10}\hypertarget{tag-005D}{}\label{cdp:eq:A.J.10}
\]
It is nonzero because the three generic tangent lines are
distinct. Equations \eqref{cdp:eq:A.J.2} and \eqref{cdp:eq:A.J.10} give
$2g(C)-2+c\leq0$, again forcing $g(C)=1$ and $c=0$
under the genus assumption. The discriminant has degree
zero and is nowhere zero. The three simple tangent roots
extend holomorphically on each small curve disc. The implicit
function theorem gives three smooth local branches of $S$
with those pairwise distinct normal tangents. Choose the
first two branches as the coordinate axes; the third is
$y=xu(x,z)$ with $u$ a unit, and replacing $y$ by $y/u(x,z)$
straightens it to $y=x$. This gives the local equation
$xy(x-y)=0$ up to a unit at every point of $D$.

\end{enumerate}
\end{proof}